% Source: https://artofproblemsolving.com/community/c3h1311481_triangle_areas
% Stable identity: aops-topic-1311481-q01
Drop the perpendicular $AH$ from $A$ to line $BC$. Because $\angle B$ is obtuse, $H$ lies on the extension of $CB$ beyond $B$, and $\angle ABH=180^\circ-120^\circ=60^\circ$.

\begin{asy}
usepackage("amsmath");
size(5.5cm);
defaultpen(fontsize(10pt));
path ra(pair P,pair u,pair v,real s=0.2){return P+s*unit(u)--P+s*unit(u)+s*unit(v)--P+s*unit(v);}
pair B=(0,0),C=(10,0),A=6*dir(120),H=(A.x,0);
draw(A--B--C--cycle,linewidth(0.9)); draw(B--H,gray+dashed); draw(A--H,red+dashed); draw(ra(H,A-H,C-H,0.4));
label("$B$",B,S); label("$C$",C,S); label("$A$",A,N); label("$H$",H,S); label("$6$",(A+B)/2,NE); label("$10$",(B+C)/2,S); label("$60^\circ$",B+(-1.1,0.5),fontsize(8pt));
label("\mbox{}",(min(currentpicture,user=true)+max(currentpicture,user=true))/2,black); // keeps the text after the figure black when LaTeX typesets the labels inline
\end{asy}

Triangle $ABH$ is a $30$-$60$-$90$ triangle with hypotenuse $AB=6$, so $BH=3$ and $AH=3\sqrt3$. Taking $BC$ as base, the height is $AH$, so
\[
[ABC]=\dfrac{1}{2}\cdot10\cdot3\sqrt3=\boxed{15\sqrt3}.
\]
